\chapter{The Campus++}

Although some may wish it, not all of your studies will take place in one building. Instead, there are a few key buildings that will play a role in your day-to-day studies.

\minisec{HSZ}
Especially the basic lectures do not take place in the APB, but in the Hörsaalzentrum (\emph{HSZ}). In this building, which is very monotonous compared to the APB, you will experience many of your lectures and accumulate one or the other kilometer while commuting between the APB and the HSZ. The HSZ has four lecture halls, including the Audimax, the largest lecture hall at the university and the largest auditorium in Saxony. In addition, there are some seminar rooms, in which tutorials can take place. Behind the \enquote{central cube} is the beautiful meadow of the HSZ\@. You can not only relax on it. Parties and fairs also take place there regularly.
In addition, the Grillcube with burgers and the like is there for small appetites, as is the Pasta-Mobile with fresh pasta.

\minisec{Willersbau}
The Willersbau~\link{https://navigator.tu-dresden.de/etplan/wil/00} is the mathematics building at the TU Dresden. This is where some of your math tutorials take place. The Willersbau is in the immediate vicinity of the HSZ and the Trefftzbau, which houses the math lecture hall.

\minisec{Trefftzbau}
This~\link{https://navigator.tu-dresden.de/etplan/tre/00} is usually where parts of your math lecture take place. Since the building is being renovated and exceptions prove the rules, no lectures are held there at the moment due to the renovation. In front of the Trefftzbau is the Trefftzwiese, which is a wonderful place to relax after a lecture, especially in the summer.

\begin{figure}[b!]
    \centering
    \includegraphics[width=\linewidth]{img/slub-lesesaal.jpg}
\end{figure}

\minisec{SLUB}
\label{sec:slub}
The Sächsische Landesbibliothek - Staats- und Universitätsbibliothek Dresden, \emph{SLUB}~\link{https://slub-dresden.de} in short, is what is simply called Bib at other universities.
With a selection of over 12 million holdings consisting of books, magazines, films, etc., it is one of the largest (university) libraries in Germany. In addition to haptic reading material, the SLUB also offers an extensive selection of online resources that you can download free of charge as a student. Many of the books that are recommended to you in basic lectures are in stock at the SLUB in large quantities. So a trip to the textbook department can save you a lot of unnecessary expenses. Across from the main building is the DrePunct, a branch of the SLUB where most computer science textbooks are stored. If you are looking for a specific book, it is advisable to check the SLUB website to find out where the book is located. This will save you most of the search and will show you whether the book you are looking for is currently available.\newline
But even for notorious non-readers, the SLUB is a popular place to stay because of the many workstations. If you want to study together with others, there is a spacious entrance area with group tables. In addition, private group rooms can be reserved. If you prefer to study alone, there are of course enough quiet places for you in the SLUB\@. The quietest workplaces can be found in the central reading room in the second basement, which can also be seen in the picture below. A nice café, where you can pay by Mensa card, completes the whole experience.

\newpage
\minisec{Seminargebäude}
The language courses of the LSK take place in the Seminargebäude, recognizable by its particularly beautiful facade. For those who want to go beyond a simple language course, additional seminars on the culture and politics of selected countries and regions are offered. More information can be found on the website of the LSK~\link{https://tu-dresden.de/gsw/slk/lsk}.
The building is right next to the SLUB.

\minisec{refectories}
Fortunately, if you don't want to spend your studies ordering pizza, you don't have to, because the Studierendenwerk operates a network of 18 refectories.
No matter in which building of the university you are, you will find a refectory or at least a café nearby.
The daily menu of all refectories can be found on the Studierendenwerk pages ~\link{https://www.studentenwerk-dresden.de/mensen/speiseplan/}.

The largest of the refectories, the Alte Mensa, is fortunately located directly between the APB and the HSZ\@. Here you will find a variety of daily changing main dishes, salads and desserts. For friends of late breakfasts or students with disturbed sleep rhythms, the Alte Mensa offers a warm evening meal until 7 p.m. between Monday and Thursday. However, the selection becomes increasingly limited after 3 p.m. Since the refectory can hardly cope with the abundance of people at peak times, it is advisable to have meal times that do \emph{not} start directly after the end of a lecture.

If you do not like what the Alte Mensa has to offer, you are welcome to take the 10-minute walk to the Zeltschlösschen. The Zeltschlösschen is a classic refectory that is inferior to the Alte Mensa in pretty much all respects.

In five minutes you can also reach the Bio-Mensa U-Boot. The ingredients used in this refectory are organic and locally produced. The meat used in the Mensa U-Boot comes directly from the west of Dresden, which increases the price of the meals considerably. However, the prices of the meatless alternatives should not differ much from those of other refectories.

You can also escape cooking on weekends, because the Mensa Siedepunkt is also open on Saturdays and Sundays. Especially after a productive study session at the SLUB, it is very easy to get there by simply changing sides of the street.

\paragraph{FSR Insider Tip:}
Without a doubt, the most popular cafeteria on campus, if you believe the statements of Dresden FSR students, is Firat.
This is an independent kebab store, but due to its efficient handling of above-average student crowds, it likes to be jokingly referred to as a refectory.
It only takes a 10-minute walk from the APB to get there and attracts with good quality, cozy ambience and a comparatively large selection of dishes.

\minisec{student club Count Down}

\includegraphics[width=\linewidth]{img/countdown.svg}

Dresden is considered the unofficial capital of student clubs; after all, there are 15 of them here.
The Studentenklub IZ e.V. also runs a small student club, the Count~Down, in the depths of the dormitory Güntzstraße 22.
IZ stands for Informatikzentrum (computer science center) and in fact this is the last remnant of the computer science faculty in Johannstadt, where it had its home until the mid-2000s.

In addition to a large selection of the best beers and the tastiest mate (drink) at student prices, the Count~Down offers a wide range of events. Various metal parties, the joint Erasmus country night with the ESN initiative of TU-Dresden as well as game and cocktail nights are part of the regular repertoire.
You can find out more about it here~\link{https://www.countdown-dresden.de/}.

Of course you should check out other student clubs as well~\link{https://vdsc.de/}.

Generally, the operation of the student clubs is ensured by voluntary commitment, so they are always looking for new blood and fresh ideas. Just get in touch with your favorite club if you can imagine being behind the bar or taking on other tasks.

\textbf{opening hours}

Mondays: 7 p.m. to midnight

Tuesdays: 8 p.m. to 1 a.m.

Wednesdays: 7 p.m. to midnight

occasionally on Saturdays
